\scanpage{358}\mspage{274}
\deleted{et pour $N$}\inserted{et tel que $f$ soit $\pi_0$-accessible}
(p.ex. le plus petit cardinal \deleted{\unsure{dit utile}}
de \unsure{type|ce type|l'espèce}). Soit alors $\pi_1$ un cardinal tel
qu'on ait $f(M_{\pi_0})\subset N_{\pi_1}$, \inserted{et $\pi_1\geq\deleted{\unsure{\sup}}\bigl(\operatorname{card\,red}M_{\pi_0}\bigr)$, \deleted{\unsure{accessible}}}
\deleted{\unsure{tel que}} \inserted{et enfin, tel que $N_{\pi_1}$ \deleted{soit une cat. \unsure{gén. finie}} soit ss-cat. \inserted{\unsure{str.|stricte}} \unsure{géné-|gén.}.}
\unsure{Alors} $(\pi_0,\pi_1)$ satisfait aux conditions énoncées.

\uline{Démonstration} \hspace{1em} On aura en effet
\[
\begin{aligned}
 M&\xleftarrow{\approx}\mathrm{Ind}_{\pi_0}(M_{\pi_0})\\
 \deleted{N}&\deleted{\xleftarrow{\approx}\mathrm{Ind}_{\pi_0}(N_{\pi_0}),}
\end{aligned}
\]
et $f$ commute\deleted{\unsure{nt|ant}} aux $\unsure{\underrightarrow{\lim}|\underline{\lim}}$ \unsure{grdes} devant
$\pi_0$, sera connu (à isom unique
près) par sa restriction à $M_{\pi_0}$, laquelle
se factorise par $N_{\pi_1}$
\[
 f\mid M_{\pi_0}: M_{\pi_0}\longrightarrow N_{\pi_1}.
\]
Pour prouver (4.12.1.1), \deleted{posant} $\pi'=\pi^{\pi_0}\geq\pi$
et notant que ${\pi'}^{\pi_0}=\pi_0$
\[
 f(M_\pi)\subset f(M_{\pi'})\overset{?}{\subset}N_{\pi'},
\]
on voit qu'il suffit de prouver
(4.12.1.2) (pour $\pi'$). \inserted{Donc \unsure{OPS} $\pi=\pi^{\pi_0}$ et prouvons (4.12.1.2).} Mais comme
$\pi_1\geq\operatorname{Card\,red}M_{\pi_0}$, on sait que si $\pi$ est
tel que $\pi\geq\pi_1$, $\pi^{\pi_0}=\pi$, alors $M_\pi$
est formé des $\underset I{\underrightarrow{\lim}}\,x_i$, les $x_i\in N_{\pi_0}$,
et $I$ \unsure{grd} devant $\unsure{\pi.|\pi_0}$ et $\operatorname{card} I\leq\pi$.
% archive sheet number: 357
